\subsection{Внедрение адаптивных функций}
\label{sec:ch1_adaptive_functions}

Описанные адаптивные функции и переход к БАВР версии 2.0 осуществлён только в терминалах производства ООО «АПС» [4]. В данной работе не представляется подробного описания, существующего функционированиях адаптивного РНМ, лишь описываются основная логика его функционирования и особые случаи наладки.

Первой базовой модернизацией было внедрение адаптивности по модулю тока, основанной на подстройке \commentnote{DISS-005-C0028}{Евдаков Алексей}{Никакой конкретики я не привожу, так как меня просили не разглашать, поэтому здесь описано только то же, что и нашем РЭ.}{уставки}. И теперь пользователь должен задавать только минимальный порог чувствительности датчиков тока, ниже которого уставка не будет опускаться. А уставка по модулю будет подстраиваться автоматически. Это позволило не только ускорить время реакции БАВР по органу РНМ, но и повысить стабильность времени функционирования, что показано в анализе, приведённом в предыдущем разделе \commentnote{DISS-005-C0029}{Евдаков Алексей}{Вставить ссылку}{(Раздел ХХ)}

\begin{figure}[htbp]
    \centering
    \includegraphics[width=0.85\textwidth]{image14.png}
    \caption[Вкладка параметрирования РНМ в терминале серии МИР]{Вкладка параметрирования РНМ в терминале \commentnote{DISS-005-C0030}{Feel}{Уставка по времени -- немного сбивает с толку. Обычно это задержка на срабатывание\newline Уставка по времени измерения? Отстройка по времени измерения?\newline Или в тексте четко термин ввести}{\commentnote{DISS-005-C0031}{Евдаков Алексей}{Лучше в тексте ввести. Ибо в ПО мы специально не совсем понятно сделали, чтобы копипастить сложнее было.}{серии}} МИР}
    \label{fig:ch1_img14}
\end{figure}

Помимо изменения амплитуды тока, так же может изменяться и характер нагрузки, например, при установке батарей статических конденсаторов (БСК) для компенсации потребляемой реактивной мощности, подстанция может перейти из режима потребления реактивной мощности, в режим генерации. А следовательно, будет изменяться оптимальный угол максимальной чувствительности (Рис.1). В большинстве же случаев угол функционирования РНМ принимает стандартное значение равное 45 градусам.

\begin{figure}[htbp]
    \centering
    \includegraphics[width=0.7\textwidth]{image15.png}
    \caption[Характеристика срабатывания РНМ]{Характеристика срабатывания РНМ}
    \label{fig:ch1_img15}
\end{figure}

Как и в случае с уставкой РНМ по амплитуде -- подстройка порога срабатывания частично может быть решена за счёт применения различных групп уставок. Однако это так же требует ручного регулирования, и вводит дополнительное влияние человеческого фактора.

\begin{figure}[htbp]
    \centering
    \includegraphics[width=0.85\textwidth]{image16.png}
    \caption[Вкладка параметрирования РНМ в терминале серии МИР]{Вкладка параметрирования РНМ в терминале серии МИР}
    \label{fig:ch1_img16}
\end{figure}

Поэтому второй базовой модернизацией было внедрение адаптивности по углу максимальной чувствительности, осуществляющей подстройку угла максимальной \commentnote{DISS-005-C0032}{Евдаков Алексей}{По этой же причине убрал какое-либо описание. Да, по хорошему оно должно быть. Но я не могу его привести.}{чувствительности}. И теперь пользователь должен задавать только базовый угол. Это позволяет повысить устойчивость функционирования БАВР.

\textbf{Примечание:} Влияние адаптивной уставки по углу альфа будет описано дальше в разборе современных проблем
